\citet{DBLP:journals/pacmpl/BernardyBNJS18} develop an extension to GHC/Haskell that
provides a notion of linear typing by using a graded type system.  The
core calculus is akin to that of Graded Base, but specialised to
the $\{0, 1, \omega\}$ semiring (Example~\ref{exm:none-one-tons}). Their calculus has no graded
modality, but the implementation as part of GHC/Haskell can
encode
the \textit{unrestricted} modality
\texttt{Ur} via an algebraic datatype:\footnote{\url{https://hackage.haskell.org/package/linear-base-0.4.0/docs/Data-Unrestricted-Linear.html\#t:Ur}}
%
\begin{haskell}
data Ur a where
  Ur :: a -> Ur a
\end{haskell}
%
where the arrow \haskin{->} denotes an $\omega$-graded arrow for the two-point semiring
$\{1, \omega\}$ ($0$ is not available in the implementation)
and thus this constructor is essentially of type $A
\xrightarrow{\omega} \mathsf{Ur}\ A$ in the Graded Base syntax. \citet{DBLP:journals/corr/abs-2112-14966} give a generalisation of this definition to a graded modality:
%
\begin{haskell}
data Box r a where
   Box :: a %r -> Box r a
\end{haskell}
%
This suggests an alternate Graded Base calculus in relation to Linear Base.
We previously showed (Section~\ref{subsec:lin-to-grad-cps}) that a type-preserving translation of Linear Base terms into
Graded Base requires a CPS encoding. By changing the substrate of Graded Base to System F, we have
\textbf{Graded~Poly~Base}, unlocking the ability for a direct-style
polymorphic Church encoding of this \texttt{Box} data type:
%
\begin{align}
\label{eq:derived-box}
   \textcolor{coeffectColor}{\square_{ \SYSTEMnt{r} } }  \SYSTEMnt{A}  =  \forall   \beta   .     (   \SYSTEMnt{A}  \xrightarrow{\textcolor{coeffectColor}{ \SYSTEMnt{r} } }   \beta    )   \xrightarrow{\textcolor{coeffectColor}{ \SYSTEMsym{1} } }   \beta    
\end{align}
%
This enables a less intrusive translation, without the need to CPS translate all terms.

\subsection{Graded Poly Base}\label{sec:graded-poly-base}
We thus extend Graded Base with polymorphic typing, the syntax and
typing of which is standard:
%
\begin{align*}
t & ::= \ldots \mid  \SYSTEMnt{t}  \text{@}  \SYSTEMnt{A}  \mid  \Lambda  \alpha  .  \SYSTEMnt{t} 
\tag{Terms} \\
\SYSTEMnt{A} & ::= \ldots \mid \alpha \mid  \forall  \alpha  .  \SYSTEMnt{B}  \tag{Types}
\end{align*}
%
where $ \alpha ,  \beta $ range over type variables, with
a corresponding standard substitution at the type-level
written $ [  \SYSTEMnt{A}  /  \alpha  ]  \SYSTEMnt{B} $.
Typing judgments are of the form $ \Delta  \vdash_{\textsc{g}_\forall }  \SYSTEMnt{t}  :  \SYSTEMnt{A} $ extending
the previous typing with:
%
\begin{align*}
\SYSTEMdruleGradPolytyAbs{} \quad
\SYSTEMdruleGradPolytyApp{}
\end{align*}
The operational semantics has the standard $\beta$ rule for type abstraction/application
and congruence:
%
\begin{align*}
\SYSTEMdruleSemGrdPolytyBeta{} \quad
\SYSTEMdruleSemGrdPolycongTyApp{}
\end{align*}
%
Introduction and elimination for the derived $\Box_r$ (from~\eqref{eq:derived-box})
can therefore be derived respectively as follows (with their typing derivations shown):
%
%\begin{align*}
%  \dfrac{ \Delta_{{\mathrm{1}}}  \vdash_{\textsc{g} }  \SYSTEMnt{t}  :  \SYSTEMnt{A} }
%        {  \textcolor{coeffectColor}{ \SYSTEMnt{r}  \cdot}  \Delta_{{\mathrm{1}}}   \vdash_{\textsc{g} }    \Lambda   \beta   .   \lambda  \SYSTEMmv{f}  .  \SYSTEMmv{f}   \,  \SYSTEMnt{t}   :    \forall   \beta   .   (   \SYSTEMnt{A}  \xrightarrow{\textcolor{coeffectColor}{ \SYSTEMnt{r} } }   \beta    )    \xrightarrow{\textcolor{coeffectColor}{ \SYSTEMsym{1} } }   \beta   }
%        \Box_i
%  \qquad
%  \dfrac{ \Delta_{{\mathrm{1}}}  \vdash_{\textsc{g} }  \SYSTEMnt{t_{{\mathrm{1}}}}  :    \forall   \beta   .   (   \SYSTEMnt{A}  \xrightarrow{\textcolor{coeffectColor}{ \SYSTEMnt{r} } }   \beta    )    \xrightarrow{\textcolor{coeffectColor}{ \SYSTEMsym{1} } }   \beta     \quad
%           \Delta_{{\mathrm{2}}} ,   \SYSTEMmv{x}  :_{\textcolor{coeffectColor}{ \SYSTEMnt{r} } }  \SYSTEMnt{A}    \vdash_{\textsc{g} }  \SYSTEMnt{t_{{\mathrm{2}}}}  :  \SYSTEMnt{B} }
%        { \Delta_{{\mathrm{1}}}  \SYSTEMsym{+}  \Delta_{{\mathrm{2}}}  \vdash_{\textsc{g} }   \SYSTEMsym{(}   \SYSTEMnt{t_{{\mathrm{1}}}}  \text{@}  \SYSTEMnt{B}   \SYSTEMsym{)} \,  \SYSTEMsym{(}   \lambda  \SYSTEMmv{x}  .  \SYSTEMnt{t_{{\mathrm{2}}}}   \SYSTEMsym{)}   :  \SYSTEMnt{B} }
%        \Box_e
%\end{align*}

\vspace{-1.5em}
\begin{align}
\inferrule*[Right=\SYSTEMRenameRuleGradPolytyAbs{}]
{
\inferrule*[Right=\SYSTEMRenameRuleGradabs{}]
{
\inferrule*[Right=\SYSTEMRenameRuleGradapp{}]
{  \Delta  \vdash_{\textsc{g}_\forall }  \SYSTEMnt{t}  :  \SYSTEMnt{A}  \quad
  \inferrule*[Right=\SYSTEMRenameRuleGradvar{}]
    {\quad}{   \SYSTEMmv{f}  :_{\textcolor{coeffectColor}{ \SYSTEMsym{1} } }    \SYSTEMnt{A}  \xrightarrow{\textcolor{coeffectColor}{ \SYSTEMnt{r} } }   \beta'       \vdash_{\textsc{g}_\forall }  \SYSTEMmv{f}  :    \SYSTEMnt{A}  \xrightarrow{\textcolor{coeffectColor}{ \SYSTEMnt{r} } }   \beta'     }}
{    \textcolor{coeffectColor}{ \SYSTEMnt{r}  \cdot}  \Delta  ,   \SYSTEMmv{f}  :_{\textcolor{coeffectColor}{ \SYSTEMsym{1} } }    \SYSTEMnt{A}  \xrightarrow{\textcolor{coeffectColor}{ \SYSTEMnt{r} } }   \beta'       \vdash_{\textsc{g}_\forall }   \SYSTEMmv{f} \,  \SYSTEMnt{t}   :   \beta'   }
}
{  \textcolor{coeffectColor}{ \SYSTEMnt{r}  \cdot}  \Delta   \vdash_{\textsc{g}_\forall }    \lambda  \SYSTEMmv{f}  .  \SYSTEMmv{f}  \,  \SYSTEMnt{t}   :    (   \SYSTEMnt{A}  \xrightarrow{\textcolor{coeffectColor}{ \SYSTEMnt{r} } }   \beta'    )   \xrightarrow{\textcolor{coeffectColor}{ \SYSTEMsym{1} } }   \beta'    }
}
{
  \textcolor{coeffectColor}{ \SYSTEMnt{r}  \cdot}  \Delta   \vdash_{\textsc{g}_\forall }    \Lambda   \beta   .   \lambda  \SYSTEMmv{f}  .  \SYSTEMmv{f}   \,  \SYSTEMnt{t}   :    \forall   \beta   .   (   \SYSTEMnt{A}  \xrightarrow{\textcolor{coeffectColor}{ \SYSTEMnt{r} } }   \beta    )    \xrightarrow{\textcolor{coeffectColor}{ \SYSTEMsym{1} } }   \beta   
}\qquad
\tag{$\Box_i$}
\end{align}
\begin{align}
\inferrule*[Right=\SYSTEMRenameRuleGradapp{}]
{
  \inferrule*[Right=\SYSTEMRenameRuleGradPolytyApp{}]
     { \Delta_{{\mathrm{1}}}  \vdash_{\textsc{g}_\forall }  \SYSTEMnt{t_{{\mathrm{1}}}}  :    \forall   \beta   .   (   \SYSTEMnt{A}  \xrightarrow{\textcolor{coeffectColor}{ \SYSTEMnt{r} } }   \beta    )    \xrightarrow{\textcolor{coeffectColor}{ \SYSTEMsym{1} } }   \beta   }
     { \Delta_{{\mathrm{1}}}  \vdash_{\textsc{g}_\forall }   \SYSTEMnt{t_{{\mathrm{1}}}}  \text{@}  \SYSTEMnt{B}   :    (   \SYSTEMnt{A}  \xrightarrow{\textcolor{coeffectColor}{ \SYSTEMnt{r} } }  \SYSTEMnt{B}   )   \xrightarrow{\textcolor{coeffectColor}{ \SYSTEMsym{1} } }  \SYSTEMnt{B}  }
  \\ \quad
  \inferrule*[Right=\SYSTEMRenameRuleGradabs{}]
    {  \Delta_{{\mathrm{2}}} ,   \SYSTEMmv{x}  :_{\textcolor{coeffectColor}{ \SYSTEMnt{r} } }  \SYSTEMnt{A}    \vdash_{\textsc{g}_\forall }  \SYSTEMnt{t_{{\mathrm{2}}}}  :  \SYSTEMnt{B} }
    { \Delta_{{\mathrm{2}}}  \vdash_{\textsc{g}_\forall }   \lambda  \SYSTEMmv{x}  .  \SYSTEMnt{t_{{\mathrm{2}}}}   :   \SYSTEMnt{A}  \xrightarrow{\textcolor{coeffectColor}{ \SYSTEMnt{r} } }  \SYSTEMnt{B}  }
}
{ \Delta_{{\mathrm{1}}}  \SYSTEMsym{+}  \Delta_{{\mathrm{2}}}  \vdash_{\textsc{g}_\forall }   \SYSTEMsym{(}   \SYSTEMnt{t_{{\mathrm{1}}}}  \text{@}  \SYSTEMnt{B}   \SYSTEMsym{)} \,  \SYSTEMsym{(}   \lambda  \SYSTEMmv{x}  .  \SYSTEMnt{t_{{\mathrm{2}}}}   \SYSTEMsym{)}   :  \SYSTEMnt{B} }
\quad
\tag{$\Box_e$}
\end{align}
%
\subsubsection{Linear Base to Graded Poly Base Translation}
\label{subsec:transl-gpb2lb}
We thus adapt the translation of types to:
%
\renewcommand{\triangleq}{& \oldTriangleq &}
\renewcommand{\equiv}{& \oldEquiv &}
\renewcommand{\SYSTEMdrule}[4][]{#3 \qquad #2}
\begin{gather*}
{\setlength{\arraycolsep}{0.1em}
\begin{array}{lcl@{\hspace{2em}}r}
\SYSTEMdruleTranslLtoGTyboxTyPoly{} & \text{(Type translation)}
\end{array}}
\end{gather*}
\let\SYSTEMdrule\oldSYSTEMdrule
\let\equiv\oldEquiv
\let\triangleq\oldTriangleq
%
We might expect that we can interpret the syntax of terms then directly
as in the above derivations, i.e.,
$ \textcolor{LGcolor}{\llparenthesis} \smidge   \textcolor{coeffectColor}{[}  \SYSTEMnt{t}  \textcolor{coeffectColor}{]}   \smidge \textcolor{LGcolor}{\rrparenthesis}  \triangleq   \Lambda   \beta   .    \lambda  \SYSTEMmv{f}  .    \SYSTEMmv{f} \,   \textcolor{LGcolor}{\llparenthesis} \smidge  \SYSTEMnt{t}  \smidge \textcolor{LGcolor}{\rrparenthesis}       $
and
$ \textcolor{LGcolor}{\llparenthesis} \smidge   \mathsf{let} \, \textcolor{coeffectColor}{[}  \SYSTEMmv{x}  \textcolor{coeffectColor}{]} =  \SYSTEMnt{t_{{\mathrm{1}}}}  \, \mathsf{in} \,  \SYSTEMnt{t_{{\mathrm{2}}}}   \smidge \textcolor{LGcolor}{\rrparenthesis}  \triangleq  \SYSTEMsym{(}    \textcolor{LGcolor}{\llparenthesis} \smidge  \SYSTEMnt{t_{{\mathrm{1}}}}  \smidge \textcolor{LGcolor}{\rrparenthesis}   \text{@}  \SYSTEMnt{B}   \SYSTEMsym{)} \,  \SYSTEMsym{(}   \lambda  \SYSTEMmv{x}  .   \textcolor{LGcolor}{\llparenthesis} \smidge  \SYSTEMnt{t_{{\mathrm{2}}}}  \smidge \textcolor{LGcolor}{\rrparenthesis}    \SYSTEMsym{)} $.
However, note the presence of $B$ on the right-hand side of $\Box$-elimination,
which has now appeared out of nowhere. Instead the translation must
ingest the entire \emph{typing derivation}, not merely the syntax of terms and types
separately. We can however translate the types separately.

Thus, we define the translation as mapping from typing derivations to terms,
where we denote derivation trees by $\Pi$:
%
\begin{align*}
\tag{Term translation}
\tikz[anchor=base, baseline]{
    \draw[line width=1pt,   LGcolor] plot[smooth, tension=.9] coordinates {(0,-0.4) (-0.1,-0.2) (-0.14,.25) (-0.1,0.7) (0,0.9)};
    \draw[line width=1pt, LGcolor] plot[smooth, tension=1] coordinates {(-.11,-0.15) (-0.13,.25) (-.11,0.65)};
    \draw[thick, LGcolor, cap=round] plot[smooth, tension=1] coordinates {(0,-0.405) (0,0.905)};
}
\dfrac
{\dfrac{\Pi}{ \Gamma  \vdash_{\textsc{l} }  \SYSTEMnt{t}  :  \SYSTEMnt{A} } \quad  \mathrm{graded}( \Gamma ) }
{  \textcolor{coeffectColor}{ \SYSTEMnt{r}  \cdot}  \Gamma   \vdash_{\textsc{l} }   \textcolor{coeffectColor}{[}  \SYSTEMnt{t}  \textcolor{coeffectColor}{]}   :   \textcolor{coeffectColor}{\square_{ \SYSTEMnt{r} } }  \SYSTEMnt{A}  }\SYSTEMRenameRuleLinpr
\tikz[anchor=base, baseline]{
    \draw[line width=1pt,   LGcolor] plot[smooth, tension=.9] coordinates {(0,-0.4) (0.1,-0.2) (0.14,.25) (0.1,0.7) (0,0.9)};
    \draw[line width=1pt, LGcolor] plot[smooth, tension=1] coordinates {(.11,-0.15) (0.13,.25) (.11,0.65)};
    \draw[thick, LGcolor, cap=round] plot[smooth, tension=1] coordinates {(0,-0.405) (0,0.905)};
}
%
& \;\triangleq\;
%
  \Lambda   \beta   .   \lambda  \SYSTEMmv{f}  .  \SYSTEMmv{f}   \,   \textcolor{LGcolor}{\llparenthesis}  \Pi  \textcolor{LGcolor}{\rrparenthesis}   \\[0.5em] %: (A r ->  tyb)
%%
\tikz[anchor=base, baseline]{
    \draw[line width=1pt,   LGcolor] plot[smooth, tension=.9] coordinates {(0,-0.4) (-0.1,-0.2) (-0.14,.25) (-0.1,0.7) (0,0.9)};
    \draw[line width=1pt, LGcolor] plot[smooth, tension=1] coordinates {(-.11,-0.15) (-0.13,.25) (-.11,0.65)};
    \draw[thick, LGcolor, cap=round] plot[smooth, tension=1] coordinates {(0,-0.405) (0,0.905)};
}
\dfrac
{\dfrac{\Pi_1}
            { \Gamma_{{\mathrm{1}}}  \vdash_{\textsc{l} }  \SYSTEMnt{t_{{\mathrm{1}}}}  :   \textcolor{coeffectColor}{\square_{ \SYSTEMnt{r} } }  \SYSTEMnt{A}  }
 \quad
 \dfrac{\Pi_2}
            {  \Gamma_{{\mathrm{2}}} ,   \SYSTEMmv{x}  : \textcolor{coeffectColor}{[}  \SYSTEMnt{A} {\textcolor{coeffectColor}{]_{ \SYSTEMnt{r} } } }    \vdash_{\textsc{l} }  \SYSTEMnt{t_{{\mathrm{2}}}}  :  \SYSTEMnt{B} }
}
{ \Gamma_{{\mathrm{1}}}  \SYSTEMsym{+}  \Gamma_{{\mathrm{2}}}  \vdash_{\textsc{l} }   \mathsf{let} \, \textcolor{coeffectColor}{[}  \SYSTEMmv{x}  \textcolor{coeffectColor}{]} =  \SYSTEMnt{t_{{\mathrm{1}}}}  \, \mathsf{in} \,  \SYSTEMnt{t_{{\mathrm{2}}}}   :  \SYSTEMnt{B} }\SYSTEMRenameRuleLinlet
\tikz[anchor=base, baseline]{
    \draw[line width=1pt,   LGcolor] plot[smooth, tension=.9] coordinates {(0,-0.4) (0.1,-0.2) (0.14,.25) (0.1,0.7) (0,0.9)};
    \draw[line width=1pt, LGcolor] plot[smooth, tension=1] coordinates {(.11,-0.15) (0.13,.25) (.11,0.65)};
    \draw[thick, LGcolor, cap=round] plot[smooth, tension=1] coordinates {(0,-0.405) (0,0.905)};
}
& \;\triangleq\;
 \SYSTEMsym{(}    \textcolor{LGcolor}{\llparenthesis}  \Pi_{{\mathrm{1}}}  \textcolor{LGcolor}{\rrparenthesis}   \text{@}   \textcolor{LGcolor}{\llparenthesis} \smidge  \SYSTEMnt{B}  \smidge \textcolor{LGcolor}{\rrparenthesis}    \SYSTEMsym{)} \,  \SYSTEMsym{(}   \lambda  \SYSTEMmv{x}  .   \textcolor{LGcolor}{\llparenthesis}  \Pi_{{\mathrm{2}}}  \textcolor{LGcolor}{\rrparenthesis}    \SYSTEMsym{)} 
\end{align*}
%
The rest of the translation follows the structure of the previous term
translation.

In the following, we write $\Pi :  \Gamma  \vdash_{\textsc{l} }  \SYSTEMnt{t}  :  \SYSTEMnt{A} $
to mean that $\Pi$ is the derivation of a judgment $ \Gamma  \vdash_{\textsc{l} }  \SYSTEMnt{t}  :  \SYSTEMnt{A} $.
%

\begin{restatable}{thrm}{linToGradTranslationPoly}[Soundness of the Linear Base to Graded Poly Base translation]
 \begin{itemize}[itemsep=0.5em,topsep=0.5em,leftmargin=1em]
  \item Type preservation: $\Pi :  \Gamma  \vdash_{\textsc{l} }  \SYSTEMnt{t}  :  \SYSTEMnt{A}  \implies   \textcolor{LGcolor}{\llparenthesis}  \Gamma  \textcolor{LGcolor}{\rrparenthesis}   \vdash_{\textsc{g}_\forall }   \textcolor{LGcolor}{\llparenthesis}  \Pi  \textcolor{LGcolor}{\rrparenthesis}   :   \textcolor{LGcolor}{\llparenthesis} \smidge  \SYSTEMnt{A}  \smidge \textcolor{LGcolor}{\rrparenthesis}  $
  \item Operational correspondence:  $(\Pi :  \Gamma  \vdash_{\textsc{l} }  \SYSTEMnt{t}  :  \SYSTEMnt{A} ) \wedge (\Pi' :  \Gamma  \vdash_{\textsc{l} }  \SYSTEMnt{t'}  :  \SYSTEMnt{A} ) \wedge ( \SYSTEMnt{t}  \rightsquigarrow_{\textsc{l} }  \SYSTEMnt{t'} ) \implies   \textcolor{LGcolor}{\llparenthesis}  \Pi  \textcolor{LGcolor}{\rrparenthesis}   \rightsquigarrow_{\textsc{g}\forall }^\ast   \textcolor{LGcolor}{\llparenthesis}  \Pi'  \textcolor{LGcolor}{\rrparenthesis}  $
  \item Equation preservation: $(\Pi :  \Gamma  \vdash_{\textsc{l} }  \SYSTEMnt{t}  :  \SYSTEMnt{A} ) \wedge (\Pi' :  \Gamma  \vdash_{\textsc{l} }  \SYSTEMnt{t'}  :  \SYSTEMnt{A} ) \wedge (\SYSTEMnt{t}  \equiv^{-\eta\Box}_{\textsc{l} }  \SYSTEMnt{t'}) \implies   \textcolor{LGcolor}{\llparenthesis}  \Pi  \textcolor{LGcolor}{\rrparenthesis}   \equiv_{\textsc{g}_\forall}   \textcolor{LGcolor}{\llparenthesis}  \Pi'  \textcolor{LGcolor}{\rrparenthesis}  $
\end{itemize}
Appendix~\refappendix{app:lin-to-grad-poly} provides the proofs.

Note the operational correspondence
needs to have typing derivations of the source term and target term.
Furthermore, the translation does not preserve $\eta$-equality for the graded modality; only the $\beta$-equality rules, commuting conversions, and congruence rules are preserved;
we thus denote the subset of the preserved equational theory here as $ \equiv^{-\eta\Box}_{\textsc{l} } $.
This is because the translation of a term $t$ of type $ \textcolor{coeffectColor}{\square_{ \SYSTEMnt{r} } }  \SYSTEMnt{A} $
may not be of the form $  \Lambda   \beta   .   \lambda  \SYSTEMmv{f}  .  \SYSTEMmv{f}   \,  \SYSTEMnt{t'} $ for some $t'$ (e.g., it could be just a free variable), which is required for $\eta$-equality to apply.
Instead, parametricity reasoning (based on the polymorphic typing) would provide that $t$ is \emph{observationally equivalent} to $  \Lambda   \beta   .   \lambda  \SYSTEMmv{f}  .  \SYSTEMmv{f}   \,  \SYSTEMnt{t'} $ for some $t'$,
and thus the translation preserves $\eta$-equality up to observational equivalence, though not up to syntactic equality;
we leave the details for further work.
\label{thrm:linToGradTranslationPoly}
\end{restatable}
%

\begin{remark}
(\haskin{Ur} is not the $!$-modality)
\citet{DBLP:journals/pacmpl/BernardyBNJS18} suggest \haskin{Ur} as
an encoding of Linear Logic's $!$-modality, although \citet{hughes2021linear} later show that this
encoding in Linear Haskell and similarly in Idris 2, as well as the graded modality in Granule, do not exactly correspond to Linear Logic's $!$-modality. Once (multiplicative) products
are included, Linear Logic does not admit distributing $!$ over products, i.e.,
$!(A \otimes B) \multimap \,! A\, \otimes \,! B$ is not derivable. However the aforementioned systems permit its derivation due to the interaction between graded modalities and products
(and similarly sums). We consider calculi with products and sums next
as this takes us closer towards practical programming and it allows us to uncover
other expressivity differences in the literature.
\end{remark}
